\documentclass[11pt,a4paper]{article}
\usepackage[utf8]{inputenc}
\usepackage[T1]{fontenc}
\usepackage[french]{babel}
\usepackage[a4paper,margin=2.3cm]{geometry}
\usepackage{graphicx}
\usepackage{xcolor}
\usepackage{booktabs}
\usepackage{enumitem}
\usepackage{titlesec}
\usepackage{fancyhdr}
\usepackage{amssymb,amsmath}

\definecolor{problue}{RGB}{20,80,150}
\definecolor{progray}{RGB}{90,90,90}
\definecolor{prolight}{RGB}{236,242,248}

\titleformat{\section}{\color{problue}\large\bfseries}{\thesection.}{0.6em}{}
\titleformat{\subsection}{\color{progray}\normalsize\bfseries}{\thesubsection}{0.6em}{}

\pagestyle{fancy}
\fancyhf{}
\lhead{\small\color{progray}ProtoVA 2 — Freinage d'urgence (AEB)}
\rhead{\small\color{progray}Douae Sebti}
\cfoot{\small\color{progray}\thepage}
\renewcommand{\headrulewidth}{0.4pt}

\begin{document}

\begin{center}
{\color{problue}\Large\bfseries Freinage d'urgence autonome (AEB) par fusion LiDAR--caméra}\\[4pt]
{\large Conception, calibration expérimentale et validation sur véhicule}\\[8pt]
{\color{progray} Douae Sebti — Projet ProtoVA 2 — séance du 24 juillet 2026}
\end{center}
\vspace{2pt}
\hrule
\vspace{10pt}

\section{Objectif et point de départ}
L'objectif de la séance était de doter le véhicule d'un \textbf{freinage d'urgence
autonome (AEB)} fiable : le véhicule doit s'arrêter seul devant tout obstacle, même
si le conducteur maintient l'accélérateur. La solution précédente d'évitement
(\emph{follow the gap}, purement réactive sur le LiDAR) s'était montrée
insuffisante : détection incomplète des obstacles et collisions avec les murs.
J'ai donc repris le problème à la base, en m'appuyant sur les capteurs embarqués
(caméra de profondeur et LiDAR) et en validant chaque choix par la mesure.

\section{Démarche : de la caméra seule à la fusion de capteurs}
\subsection{Détection générique par la profondeur}
La détection YOLO seule ne convient pas pour un AEB : elle ne reconnaît que les
classes apprises (piéton, chaise\ldots) et \textbf{ignore les murs}. J'ai donc ajouté
au nœud de perception une mesure directe sur l'\textbf{image de profondeur} : la
distance au plus proche obstacle dans un couloir frontal est extraite de la carte de
profondeur brute (percentile bas, zone centrale), indépendamment de toute
classification. Tout obstacle, connu ou non, est ainsi détecté.

\subsection{Limite de la caméra et fusion avec le LiDAR}
Les essais ont montré une limite physique de la caméra temps-de-vol : les surfaces
absorbant l'infrarouge (caoutchouc, matériaux sombres) renvoient une profondeur
invalide et deviennent invisibles. J'ai donc conçu un nœud d'AEB par
\textbf{fusion de capteurs}, exécuté localement sur la Jetson (réaction $\sim$30~Hz,
sans dépendre du PC ni du réseau) :
\begin{itemize}[leftmargin=1.4em,itemsep=1pt]
  \item \textbf{LiDAR} (balayage 360°) : distance au plus proche obstacle dans les
        secteurs avant \emph{et} arrière ($\pm$45°) — fiable sur toutes les surfaces ;
  \item \textbf{caméra de profondeur} : couloir frontal, en complément (objets hors
        du plan du LiDAR) ;
  \item distance de danger retenue $=$ \textbf{minimum des deux capteurs}, par direction.
\end{itemize}
Un point matériel important a été identifié au passage : le LiDAR étant monté
« moteur à l'arrière », l'avant réel du véhicule correspond à \textbf{180°} du repère
du scan. Je l'ai vérifié sur données réelles (obstacle à 0,67~m derrière contre
3,2~m d'espace libre devant) et corrigé dans tous les nœuds concernés.

\section{Calibration du freinage au banc d'essai}
Premier constat en essai : couper les gaz ne freine pas — le variateur (ESC) passe
en \textbf{roue libre} et le véhicule continue sur son élan. J'ai donc mis en place un
\textbf{banc de mesure} (roues levées, chronométrage par l'encodeur de roue) pour
caractériser l'arrêt et calibrer un \textbf{freinage actif} par contre-impulsion :

\begin{center}
\renewcommand{\arraystretch}{1.2}
\begin{tabular}{@{}lcc@{}}
\toprule
\textbf{Stratégie de freinage} & \textbf{Temps d'arrêt} & \textbf{Gain} \\
\midrule
Roue libre (gaz coupés)              & 1,94 s & --- \\
Contre-impulsion $-0{,}3$ / 0,25 s   & 1,16 s & $-40\,\%$ \\
\textbf{Contre-impulsion $-0{,}5$ / 0,5 s} & \textbf{0,61 s} & $\mathbf{-69\,\%}$ \\
Contre-impulsion $-0{,}8$ / 0,5 s    & 0,57 s & saturation ESC \\
\bottomrule
\end{tabular}
\end{center}

Le réglage $-0{,}5$ pendant 0,5~s est retenu (au-delà, l'ESC sature). L'impulsion est
\textbf{interrompue dès l'arrêt réel} mesuré par l'encodeur, et n'est appliquée qu'en
marche avant : en freinage de marche arrière, l'arrêt est direct, sans aucun
mouvement inverse parasite.

\section{Comportement final de l'AEB}
\begin{itemize}[leftmargin=1.4em,itemsep=2pt]
  \item \textbf{Bidirectionnel} : la marche avant est freinée par les obstacles
        détectés devant, la marche arrière par ceux détectés derrière (LiDAR 360°).
        Les deux distances sont affichées en continu à l'écran.
  \item \textbf{Prioritaire} : la commande d'arrêt passe par le multiplexeur de
        commandes (\texttt{twist\_mux}) avec la priorité maximale (255) — le véhicule
        \textbf{reste arrêté même accélérateur enfoncé}.
  \item \textbf{Échappatoire} : la direction opposée au danger reste toujours
        disponible pour s'écarter de l'obstacle.
  \item \textbf{Distance de freinage dynamique} : les seuils croissent avec la
        vitesse réelle mesurée par l'encodeur
        ($d_{\text{arrêt}} = d_0 + k\,|v|$) — plus le véhicule roule vite, plus il
        freine tôt, comme un AEB réel.
  \item \textbf{Réactivité} : j'ai également réduit le délai de reprise en main après
        un arrêt d'urgence (temporisation du multiplexeur ramenée de 1~s à 0,3~s).
\end{itemize}

\section{Validation par enregistrement embarqué}
Pour trancher entre ressenti de conduite et comportement réel, j'ai enregistré
pendant la conduite une « boîte noire » synchronisant les distances LiDAR, la
commande du conducteur, la sortie de l'AEB, la commande moteur et la vitesse
encodeur. L'enregistrement confirme le fonctionnement :

\begin{center}
\fcolorbox{problue}{prolight}{\parbox{0.93\textwidth}{\ttfamily\small
t=47,1 s \ avant 0,77 m \ pédale +0,16 \ AEB $-$0,50 (impulsion de frein) \ roues 0,35\\
t=48,2 s \ avant 0,66 m \ \textbf{pédale +0,20 (enfoncée)} \ AEB 0,00 \ \textbf{roues 0,00 — arrêté}
}}
\end{center}
\vspace{2pt}
\begin{center}\small\color{progray}
Extrait : arrêt complet devant l'obstacle alors que l'accélérateur reste enfoncé.
\end{center}

Cette méthode a aussi révélé le délai de redémarrage d'une seconde (temporisation du
multiplexeur), corrigé ensuite.

\section{Bilan et suites}
\textbf{Bilan.} Le véhicule dispose désormais d'un freinage d'urgence autonome
bidirectionnel, fondé sur la fusion LiDAR--caméra, calibré expérimentalement
(temps d'arrêt réduit de 1,94~s à 0,61~s) et validé en conduite réelle par
enregistrement embarqué. L'ensemble du code et des mesures est archivé sur le dépôt
Git du projet.

\textbf{Suites envisagées.}
\begin{itemize}[leftmargin=1.4em,itemsep=1pt]
  \item affiner les coefficients de seuil dynamique ($k$) sur davantage d'essais au sol ;
  \item vérifier le comportement du servomoteur de direction (point resté ouvert) ;
  \item intégrer l'AEB à la démonstration V2V complète (freinage $\rightarrow$ message
        DENM vers le véhicule voisin).
\end{itemize}

\end{document}
